\documentclass[12pt,a4paper]{article}
\usepackage[utf8]{inputenc}
\usepackage[vietnamese]{babel}
\usepackage[left=2cm,right=2cm,top=2cm,bottom=2cm]{geometry}
\usepackage{mathptmx}
\usepackage{enumitem}

\pagestyle{plain}
\linespread{1.15}

\begin{document}

% --- HEADER QUỐC HIỆU ---
\noindent
\begin{minipage}[t]{0.45\textwidth}
    \begin{center}
        \small
        BỘ CÔNG AN\\
        \textbf{\underline{VIỆN KHOA HỌC HÌNH SỰ}}\\[0.1cm]
        Số: 09/KL-KTHS\\
        \textit{V/v kết luận giám định tài liệu}
    \end{center}
\end{minipage}
\hfill
\begin{minipage}[t]{0.52\textwidth}
    \begin{center}
        \small
        \textbf{CỘNG HÒA XÃ HỘI CHỦ NGHĨA VIỆT NAM}\\
        \textbf{\underline{Độc lập - Tự do - Hạnh phúc}}\\[0.1cm]
        \textit{Hà Nội, ngày 25 tháng 07 năm 2026}
    \end{center}
\end{minipage}

\vspace{0.6cm}

\begin{center}
    \textbf{\large KẾT LUẬN GIÁM ĐỊNH KỸ THUẬT HÌNH SỰ}\\
    \textit{(V/v Giám định tính thật giả và thời điểm tạo lập tài liệu ký hiệu f2-1)}
\end{center}

\vspace{0.3cm}

\noindent Viện Khoa học Hình sự Bộ Công an kết luận về mẫu giám định di chúc gửi kèm Quyết định trưng cầu số 14/QĐ-CSHS như sau:

\section*{I. NỘI DUNG YÊU CẦU GIÁM ĐỊNH}
\begin{enumerate}[label=\arabic*., leftmargin=1.5cm]
    \item Chữ viết, chữ ký trên tài liệu di chúc ký hiệu \textbf{A} (lập ngày 15/04/2018) có bị tẩy xóa, sửa chữa hoặc chèn đè nội dung hay không?
    \item Loại mực và thời điểm tạo lập đoạn văn bản *"toàn quyền sở hữu và hưởng thụ tiền đền bù"* ở dòng thứ 8 có đồng nhất với phần chữ viết gốc trên tài liệu không?
\end{enumerate}

\section*{II. KẾT QUẢ GIÁM ĐỊNH}
\begin{enumerate}[label=\arabic*., leftmargin=1.5cm]
    \item \textbf{Khảo sát bề mặt tài liệu A dưới kính phổ quang VSC/UV:} Phát hiện diện tích tổn thương xơ sợi giấy kích thước $1,5 \times 4,2\text{ cm}$ tại dòng thứ 8. Bề mặt giấy tồn dư hợp chất Bleaching Agent (chất tẩy rửa hóa học gốc Clo chuyên dụng).
    \item \textbf{Phân tích sắc ký lớp mỏng (TLC) \& Độ lão hóa mực:}
    \begin{itemize}[leftmargin=0.5cm]
        \item Nét chữ gốc của cụ Nguyễn Văn Thành được viết bằng mực bút máy bơm gốc sắt-gallo (Iron gall ink), có độ lão hóa phù hợp thời điểm năm 2018 (đồng nhất mẫu so sánh \textbf{M1}).
        \item Đoạn văn bản bổ sung tại dòng thứ 8 được viết bằng mực bi hóa dầu (Oil-based ballpoint ink) mang đặc tính quang học sản xuất từ năm 2024, viết chèn đè lên vùng giấy đã bị tẩy xóa hóa chất.
    \end{itemize}
\end{enumerate}

\section*{III. KẾT LUẬN GIÁM ĐỊNH}
\begin{enumerate}[label=\arabic*., leftmargin=1.5cm]
    \item Tài liệu di chúc ký hiệu \textbf{A} đã bị \textbf{TẨY XÓA HÓA CHẤT} tại vị trí dòng thứ 8 (khu vực ghi thông tin người thừa kế ban đầu).
    \item Đoạn văn bản *"toàn quyền sở hữu và hưởng thụ tiền đền bù"* tại dòng thứ 8 được \textbf{VIẾT CHÈN ĐÈ BỔ SUNG} sau thời điểm lập di chúc gốc, sử dụng mực bi sản xuất từ năm 2024.
\end{enumerate}

\vspace{0.8cm}

\noindent
\begin{minipage}[t]{0.45\textwidth}
    \small
    \textbf{\textit{Nơi nhận:}}\\
    - Cơ quan CSĐT Công an TP;\\
    - Lưu: VT, CSHS.
\end{minipage}
\hfill
\begin{minipage}[t]{0.5\textwidth}
    \begin{center}
        \small
        \textbf{VIỆN TRƯỞNG VIỆN KHHS}\\
        \textit{(Ký và đóng dấu)}\\[1.5cm]
        \textbf{Thiếu tướng Trần Quốc Tuấn}
    \end{center}
\end{minipage}

\end{document}
